\documentclass[11pt]{article}
\usepackage{amsmath}
\usepackage{amssymb}
\usepackage{geometry}
\usepackage{graphicx}
\usepackage{xcolor}
\usepackage{booktabs}

\geometry{a4paper, margin=1in}

\title{\textbf{Homework 4.4: Coordinate Systems}}
\author{Tashfeen Omran}
\date{September 2026}

\begin{document}

\maketitle
\tableofcontents
\newpage

\section{Part 1: Comprehensive Introduction, Context, and Prerequisites}

\subsection{Core Concepts}
In linear algebra, a \textbf{basis} $B$ for a vector space $V$ provides a frame of reference—a "coordinate system"—for that space. The fundamental idea of Section 4.4 is that once a basis is chosen, every vector in that abstract space can be uniquely represented by a column of numbers (its coordinates). 

If $B = \{b_1, b_2, \dots, b_n\}$ is a basis for a vector space $V$, the \textbf{Unique Representation Theorem} states that for every vector $x$ in $V$, there exists a \textit{unique} set of scalars $c_1, c_2, \dots, c_n$ such that:
\[ x = c_1 b_1 + c_2 b_2 + \dots + c_n b_n \]
These scalars are called the \textbf{$B$-coordinates of $x$}. We package these scalars into a vector in $\mathbb{R}^n$, denoted as $[x]_B$:
\[ [x]_B = \begin{bmatrix} c_1 \\ c_2 \\ \vdots \\ c_n \end{bmatrix} \]

\subsection{Intuition and Derivation}
Why do we need this? Imagine you are looking at a plane in $\mathbb{R}^3$. Vectors in this plane have 3 components (e.g., $x, y, z$), but the plane itself is only 2-dimensional. By choosing a basis of two vectors for this plane, we can describe any point on the plane using just 2 coordinates instead of 3. 

The matrix that connects our new coordinate system to the standard coordinate system is called the \textbf{change-of-coordinates matrix}, denoted $P_B$. It is formed by placing the basis vectors of $B$ as its columns:
\[ P_B = \begin{bmatrix} b_1 & b_2 & \dots & b_n \end{bmatrix} \]
The derivation of the coordinate transformation is straightforward matrix-vector multiplication. By the definition of matrix multiplication:
\[ P_B [x]_B = \begin{bmatrix} b_1 & b_2 & \dots & b_n \end{bmatrix} \begin{bmatrix} c_1 \\ c_2 \\ \vdots \\ c_n \end{bmatrix} = c_1 b_1 + c_2 b_2 + \dots + c_n b_n = x \]
Thus, we get the master equation: $x = P_B [x]_B$. Since the columns of $P_B$ form a basis, they are linearly independent, meaning $P_B$ is invertible. Multiplying both sides by $P_B^{-1}$ gives us the inverse mapping: $[x]_B = P_B^{-1} x$.

\subsection{Historical Context and Motivation}
The concept of coordinates dates back to René Descartes in the 17th century, who bridged algebra and geometry by introducing the Cartesian coordinate system. However, in the 19th century, mathematicians like Hermann Grassmann and Giuseppe Peano began studying abstract vector spaces—spaces of polynomials, matrices, or continuous functions. They needed a way to perform concrete calculations on these abstract objects. The coordinate mapping $x \mapsto [x]_B$ was developed as a bridge. It acts as an \textbf{isomorphism} (a structure-preserving bijection), allowing mathematicians to translate an abstract problem into a standard $\mathbb{R}^n$ matrix problem, solve it there, and translate the answer back.

\subsection{Key Formulas}
\begin{itemize}
    \item \textbf{Vector Reconstruction:} $x = P_B [x]_B$
    \item \textbf{Coordinate Extraction:} $[x]_B = P_B^{-1} x$
    \item \textbf{Change-of-Coordinates Matrix:} $P_B = \begin{bmatrix} b_1 & b_2 & \dots & b_n \end{bmatrix}$
    \item \textbf{$2 \times 2$ Matrix Inverse:} If $A = \begin{bmatrix} a & b \\ c & d \end{bmatrix}$, then $A^{-1} = \frac{1}{ad-bc} \begin{bmatrix} d & -b \\ -c & a \end{bmatrix}$
\end{itemize}

\subsection{Prerequisites}
To succeed in this chapter, you must be perfectly comfortable with:
1. \textbf{Row Reduction (Gaussian Elimination):} Solving augmented matrices $[A \mid b]$.
2. \textbf{Matrix Inversion:} Especially the $2 \times 2$ shortcut formula.
3. \textbf{Linear Independence and Span:} Understanding what makes a set of vectors a basis.
4. \textbf{Polynomial Algebra:} Equating coefficients of like terms.

\newpage
\section{Part 2: Detailed Homework Solutions}

\textit{Note: A strategic selection of problems has been chosen to cover every distinct archetype present in the homework set, ensuring comprehensive mastery without redundant repetition.}

\subsection{Problem 2: Finding the Vector $x$ from its Coordinates}
\textbf{Problem Statement:} Let $B = \{b_1, b_2\}$ where $b_1 = \begin{bmatrix} 4 \\ 5 \end{bmatrix}$ and $b_2 = \begin{bmatrix} 6 \\ 7 \end{bmatrix}$. Find the vector $x$ determined by the given coordinate vector $[x]_B = \begin{bmatrix} 8 \\ -5 \end{bmatrix}$.

\textbf{Solution:}
By the definition of coordinate vectors, the entries of $[x]_B$ are the weights $c_1, c_2$ in the linear combination of the basis vectors.
\[ x = c_1 b_1 + c_2 b_2 \]
Substitute the given values:
\[ x = 8 \begin{bmatrix} 4 \\ 5 \end{bmatrix} + (-5) \begin{bmatrix} 6 \\ 7 \end{bmatrix} \]
Distribute the scalars:
\[ x = \begin{bmatrix} 32 \\ 40 \end{bmatrix} + \begin{bmatrix} -30 \\ -35 \end{bmatrix} \]
Add the vectors:
\[ x = \begin{bmatrix} 32 - 30 \\ 40 - 35 \end{bmatrix} = \begin{bmatrix} 2 \\ 5 \end{bmatrix} \]
\textbf{Final Answer:} $x = \begin{bmatrix} 2 \\ 5 \end{bmatrix}$

\subsection{Problem 5: Finding Coordinates via Row Reduction}
\textbf{Problem Statement:} Find the coordinate vector $[x]_B$ of $x$ relative to the given basis $B = \{b_1, b_2\}$, where $b_1 = \begin{bmatrix} 1 \\ -3 \end{bmatrix}$, $b_2 = \begin{bmatrix} 2 \\ -5 \end{bmatrix}$, and $x = \begin{bmatrix} -2 \\ 1 \end{bmatrix}$.

\textbf{Solution:}
We need to find scalars $c_1, c_2$ such that $c_1 b_1 + c_2 b_2 = x$. This is equivalent to solving the matrix equation $P_B [x]_B = x$. We set up the augmented matrix $[b_1 \ b_2 \mid x]$:
\[ \begin{bmatrix} 1 & 2 & | & -2 \\ -3 & -5 & | & 1 \end{bmatrix} \]
Perform row operations to reach reduced row echelon form (RREF).
$R_2 \leftarrow R_2 + 3R_1$:
\[ \begin{bmatrix} 1 & 2 & | & -2 \\ 0 & 1 & | & -5 \end{bmatrix} \]
$R_1 \leftarrow R_1 - 2R_2$:
\[ \begin{bmatrix} 1 & 0 & | & 8 \\ 0 & 1 & | & -5 \end{bmatrix} \]
The rightmost column gives the coordinates.
\textbf{Final Answer:} $[x]_B = \begin{bmatrix} 8 \\ -5 \end{bmatrix}$

\subsection{Problem 9: The Change-of-Coordinates Matrix}
\textbf{Problem Statement:} Find the change-of-coordinates matrix from $B$ to the standard basis in $\mathbb{R}^2$, given $B = \left\{ \begin{bmatrix} 2 \\ -9 \end{bmatrix}, \begin{bmatrix} 1 \\ 8 \end{bmatrix} \right\}$.

\textbf{Solution:}
The change-of-coordinates matrix $P_B$ is constructed by simply placing the basis vectors of $B$ as the columns of the matrix.
\[ P_B = \begin{bmatrix} b_1 & b_2 \end{bmatrix} = \begin{bmatrix} 2 & 1 \\ -9 & 8 \end{bmatrix} \]
\textbf{Final Answer:} $P_B = \begin{bmatrix} 2 & 1 \\ -9 & 8 \end{bmatrix}$

\subsection{Problem 11: Finding Coordinates via Matrix Inverse}
\textbf{Problem Statement:} Use an inverse matrix to find $[x]_B$ for $B = \left\{ \begin{bmatrix} 3 \\ -5 \end{bmatrix}, \begin{bmatrix} -4 \\ 6 \end{bmatrix} \right\}$ and $x = \begin{bmatrix} 2 \\ -6 \end{bmatrix}$.

\textbf{Solution:}
First, construct the change-of-coordinates matrix $P_B$:
\[ P_B = \begin{bmatrix} 3 & -4 \\ -5 & 6 \end{bmatrix} \]
We know that $[x]_B = P_B^{-1} x$. Let's find $P_B^{-1}$ using the $2 \times 2$ inverse formula:
\[ \det(P_B) = (3)(6) - (-4)(-5) = 18 - 20 = -2 \]
\[ P_B^{-1} = \frac{1}{-2} \begin{bmatrix} 6 & 4 \\ 5 & 3 \end{bmatrix} = \begin{bmatrix} -3 & -2 \\ -2.5 & -1.5 \end{bmatrix} \]
Now, multiply $P_B^{-1}$ by $x$:
\[ [x]_B = \begin{bmatrix} -3 & -2 \\ -2.5 & -1.5 \end{bmatrix} \begin{bmatrix} 2 \\ -6 \end{bmatrix} \]
\[ [x]_B = \begin{bmatrix} (-3)(2) + (-2)(-6) \\ (-2.5)(2) + (-1.5)(-6) \end{bmatrix} = \begin{bmatrix} -6 + 12 \\ -5 + 9 \end{bmatrix} = \begin{bmatrix} 6 \\ 4 \end{bmatrix} \]
\textbf{Final Answer:} $[x]_B = \begin{bmatrix} 6 \\ 4 \end{bmatrix}$

\subsection{Problem 13: Coordinates in Abstract Polynomial Spaces}
\textbf{Problem Statement:} The set $B = \{1+t^2, t+t^2, 1+2t+t^2\}$ is a basis for $\mathbb{P}_2$. Find the coordinate vector of $p(t) = 1+4t+7t^2$ relative to $B$.

\textbf{Solution:}
We must find scalars $c_1, c_2, c_3$ such that:
\[ c_1(1+t^2) + c_2(t+t^2) + c_3(1+2t+t^2) = 1+4t+7t^2 \]
Expand the left side and group by powers of $t$:
\[ (c_1 + c_3) + (c_2 + 2c_3)t + (c_1 + c_2 + c_3)t^2 = 1 + 4t + 7t^2 \]
Equating the coefficients of corresponding powers of $t$ yields a system of linear equations:
\begin{align*}
c_1 + c_3 &= 1 \quad \text{(constant term)} \\
c_2 + 2c_3 &= 4 \quad \text{($t$ term)} \\
c_1 + c_2 + c_3 &= 7 \quad \text{($t^2$ term)}
\end{align*}
Set up the augmented matrix and row reduce:
\[ \begin{bmatrix} 1 & 0 & 1 & | & 1 \\ 0 & 1 & 2 & | & 4 \\ 1 & 1 & 1 & | & 7 \end{bmatrix} \]
$R_3 \leftarrow R_3 - R_1$:
\[ \begin{bmatrix} 1 & 0 & 1 & | & 1 \\ 0 & 1 & 2 & | & 4 \\ 0 & 1 & 0 & | & 6 \end{bmatrix} \]
$R_3 \leftarrow R_3 - R_2$:
\[ \begin{bmatrix} 1 & 0 & 1 & | & 1 \\ 0 & 1 & 2 & | & 4 \\ 0 & 0 & -2 & | & 2 \end{bmatrix} \]
$R_3 \leftarrow R_3 / -2$:
\[ \begin{bmatrix} 1 & 0 & 1 & | & 1 \\ 0 & 1 & 2 & | & 4 \\ 0 & 0 & 1 & | & -1 \end{bmatrix} \]
$R_2 \leftarrow R_2 - 2R_3$:
\[ \begin{bmatrix} 1 & 0 & 1 & | & 1 \\ 0 & 1 & 0 & | & 6 \\ 0 & 0 & 1 & | & -1 \end{bmatrix} \]
$R_1 \leftarrow R_1 - R_3$:
\[ \begin{bmatrix} 1 & 0 & 0 & | & 2 \\ 0 & 1 & 0 & | & 6 \\ 0 & 0 & 1 & | & -1 \end{bmatrix} \]
The solution is $c_1 = 2, c_2 = 6, c_3 = -1$.
\textbf{Final Answer:} $[p]_B = \begin{bmatrix} 2 \\ 6 \\ -1 \end{bmatrix}$

\subsection{Problem 15 \& 16: Conceptual True/False}
\textbf{Problem 15 Statement:} (T/F) If $x$ is in $V$ and if $B$ contains $n$ vectors, then the $B$-coordinate vector of $x$ is in $\mathbb{R}^n$.
\textbf{Solution:} \textbf{True.} By the Unique Representation Theorem, $x$ is written as a linear combination of the $n$ basis vectors. The $n$ scalar weights form a column vector with $n$ entries, which by definition is an element of $\mathbb{R}^n$.

\textbf{Problem 16 Statement:} (T/F) If $B$ is the standard basis for $\mathbb{R}^n$, then the $B$-coordinate vector of an $x$ in $\mathbb{R}^n$ is $x$ itself.
\textbf{Solution:} \textbf{True.} The standard basis $E = \{e_1, \dots, e_n\}$ forms the identity matrix $P_E = I$. Since $[x]_E = P_E^{-1} x = I^{-1} x = Ix = x$, the coordinate vector is identical to the vector itself.

\subsection{Problem 27: Proof of One-to-One Mapping}
\textbf{Problem Statement:} Show that the coordinate mapping is one-to-one. (Hint: Suppose $[u]_B = [w]_B$ for some $u$ and $w$ in $V$, and show that $u = w$.)

\textbf{Solution:}
Let $B = \{b_1, b_2, \dots, b_n\}$ be a basis for a vector space $V$. 
Assume that for two vectors $u, w \in V$, their coordinate vectors are equal:
\[ [u]_B = [w]_B = \begin{bmatrix} c_1 \\ c_2 \\ \vdots \\ c_n \end{bmatrix} \]
By the definition of the coordinate mapping, the vector $u$ is constructed using these weights:
\[ u = c_1 b_1 + c_2 b_2 + \dots + c_n b_n \]
Similarly, the vector $w$ is constructed using the exact same weights:
\[ w = c_1 b_1 + c_2 b_2 + \dots + c_n b_n \]
Since the right-hand sides of both equations are identical, it strictly follows that $u = w$. Because $f(u) = f(w) \implies u = w$, the coordinate mapping is one-to-one.

\subsection{Problem 31: Linear Independence via Isomorphism}
\textbf{Problem Statement:} Use coordinate vectors to test whether the following set of polynomials is linearly independent: $1+2t^3, 2+t-3t^2, -t+2t^2-t^3$.

\textbf{Solution:}
We map these polynomials to $\mathbb{R}^4$ using the standard basis $E = \{1, t, t^2, t^3\}$ for $\mathbb{P}_3$.
\begin{align*}
p_1(t) &= 1 + 0t + 0t^2 + 2t^3 \implies [p_1]_E = \begin{bmatrix} 1 \\ 0 \\ 0 \\ 2 \end{bmatrix} \\
p_2(t) &= 2 + 1t - 3t^2 + 0t^3 \implies [p_2]_E = \begin{bmatrix} 2 \\ 1 \\ -3 \\ 0 \end{bmatrix} \\
p_3(t) &= 0 - 1t + 2t^2 - 1t^3 \implies [p_3]_E = \begin{bmatrix} 0 \\ -1 \\ 2 \\ -1 \end{bmatrix}
\end{align*}
To check for linear independence, we place these coordinate vectors as columns in a matrix $A$ and row reduce to see if every column has a pivot.
\[ A = \begin{bmatrix} 1 & 2 & 0 \\ 0 & 1 & -1 \\ 0 & -3 & 2 \\ 2 & 0 & -1 \end{bmatrix} \]
$R_4 \leftarrow R_4 - 2R_1$:
\[ \begin{bmatrix} 1 & 2 & 0 \\ 0 & 1 & -1 \\ 0 & -3 & 2 \\ 0 & -4 & -1 \end{bmatrix} \]
$R_3 \leftarrow R_3 + 3R_2$ and $R_4 \leftarrow R_4 + 4R_2$:
\[ \begin{bmatrix} 1 & 2 & 0 \\ 0 & 1 & -1 \\ 0 & 0 & -1 \\ 0 & 0 & -5 \end{bmatrix} \]
$R_4 \leftarrow R_4 - 5R_3$:
\[ \begin{bmatrix} 1 & 2 & 0 \\ 0 & 1 & -1 \\ 0 & 0 & -1 \\ 0 & 0 & 0 \end{bmatrix} \]
The matrix has a pivot in every column (columns 1, 2, and 3). Therefore, the coordinate vectors are linearly independent in $\mathbb{R}^4$. Because the coordinate mapping is an isomorphism (it preserves linear relations), the original polynomials are linearly independent in $\mathbb{P}_3$.
\textbf{Final Answer:} The polynomials are linearly independent.

\newpage
\section{Part 3: In-Depth Analysis of Problems and Techniques}

\subsection{A) Problem Types and General Approach}
\begin{itemize}
    \item \textbf{Type 1: Vector Reconstruction (e.g., Problem 2).} 
    \textit{Strategy:} You are given the basis $B$ and the coordinates $[x]_B$. Simply multiply the change-of-coordinates matrix $P_B$ by $[x]_B$. This is a direct application of the definition $x = c_1 b_1 + \dots + c_n b_n$.
    
    \item \textbf{Type 2: Coordinate Extraction via Row Reduction (e.g., Problems 5, 13).} 
    \textit{Strategy:} You are given $x$ and $B$, and need to find $[x]_B$. Set up the augmented matrix $[P_B \mid x]$ and row reduce to $[I \mid [x]_B]$. This is the most robust method and works for any dimension.
    
    \item \textbf{Type 3: Coordinate Extraction via Matrix Inverse (e.g., Problem 11).} 
    \textit{Strategy:} Use the formula $[x]_B = P_B^{-1} x$. This is highly efficient for $2 \times 2$ matrices because of the explicit inverse formula, but computationally expensive for larger matrices compared to row reduction.
    
    \item \textbf{Type 4: Isomorphic Mapping for Abstract Spaces (e.g., Problems 13, 31).} 
    \textit{Strategy:} Translate polynomials (or matrices) into column vectors in $\mathbb{R}^n$ by extracting their coefficients relative to a standard basis. Solve the problem in $\mathbb{R}^n$, then translate the conclusion back to the abstract space.
\end{itemize}

\subsection{B) Key Algebraic and Calculus Manipulations}
\begin{itemize}
    \item \textbf{Equating Coefficients (Problem 13):} When dealing with polynomials, two polynomials are equal if and only if their corresponding coefficients are equal. This algebraic truth allows us to convert a single polynomial equation into a system of linear equations.
    \item \textbf{The $2 \times 2$ Inverse Shortcut (Problem 11):} Swapping the main diagonal, negating the anti-diagonal, and dividing by the determinant is a crucial time-saver.
    \item \textbf{Pivot Analysis (Problem 31):} To prove linear independence, we don't need to fully solve the system; we only need to row reduce far enough to confirm that every column contains a pivot position.
\end{itemize}

\newpage
\section{Part 4: "Cheatsheet" and Tips for Success}

\subsection{Summary of Important Formulas}
\begin{table}[h]
\centering
\begin{tabular}{ll}
\toprule
\textbf{Concept} & \textbf{Formula} \\
\midrule
Change-of-Coordinates Matrix & $P_B = \begin{bmatrix} b_1 & b_2 & \dots & b_n \end{bmatrix}$ \\
Vector from Coordinates & $x = P_B [x]_B$ \\
Coordinates from Vector & $[x]_B = P_B^{-1} x$ \\
Augmented Matrix Setup & $[P_B \mid x] \sim [I \mid [x]_B]$ \\
\bottomrule
\end{tabular}
\end{table}

\subsection{Tricks and Shortcuts}
\begin{itemize}
    \item \textbf{The $2 \times 2$ Rule:} If asked to find $[x]_B$ for a 2-dimensional space, immediately use the $P_B^{-1}$ formula. It is faster and less prone to arithmetic errors than row reduction.
    \item \textbf{Polynomial Translation:} When mapping polynomials to vectors, always write out the "missing" terms with a $0$ coefficient (e.g., $1+t^2 \to 1 + 0t + 1t^2 \to [1, 0, 1]^T$). This prevents shifting errors.
\end{itemize}

\subsection{Common Pitfalls to Watch Out For}
\begin{itemize}
    \item \textbf{Reversing the Mapping:} The most common mistake is calculating $P_B x$ when you are asked to find $[x]_B$. Remember: $P_B$ takes you \textit{from} the $B$-basis \textit{to} the standard basis. $P_B^{-1}$ takes you \textit{from} the standard basis \textit{to} the $B$-basis.
    \item \textbf{Basis Order Matters:} A basis is an \textit{ordered} set. If you swap the order of $b_1$ and $b_2$, the entries in $[x]_B$ will also swap. Always respect the given order.
\end{itemize}

\newpage
\section{Part 5: Conceptual Synthesis and The "Big Picture"}

\subsection{A) Thematic Connections}
The overarching theme of this section is \textbf{Isomorphism}—the idea that two mathematically distinct objects can have the exact same structural properties. By proving that an $n$-dimensional vector space $V$ is isomorphic to $\mathbb{R}^n$, we establish that $\mathbb{R}^n$ is the "universal template" for all finite-dimensional vector spaces. Whether you are studying quantum states, continuous functions, or geometric planes, you can map them to $\mathbb{R}^n$, use standard matrix operations, and map the results back. This is analogous to using a logarithm to turn a difficult multiplication problem into a simple addition problem, and then exponentiating to get the final answer.

\subsection{B) Forward and Backward Links}
\begin{itemize}
    \item \textbf{Backward Links:} This topic relies heavily on the definitions of \textit{linear independence} and \textit{span}. If the columns of $P_B$ did not span the space, we couldn't represent every vector. If they weren't linearly independent, the representation wouldn't be unique (and $P_B$ wouldn't be invertible).
    \item \textbf{Forward Links:} This is the bedrock for \textbf{Eigenvectors and Diagonalization} (Chapter 5). Soon, we will actively search for a "magic" basis $B$ (an eigenbasis) that makes a complex linear transformation look like a simple diagonal matrix. The formula $A = P D P^{-1}$ is exactly a change of coordinates into the eigenbasis, a scaling operation ($D$), and a change back to the standard basis.
\end{itemize}

\newpage
\section{Part 6: Abstract Real-World Modeling (AI, ML, \& Quant Finance)}

\subsection{A) Concrete Scenarios}
\begin{enumerate}
    \item \textbf{Machine Learning (Principal Component Analysis - PCA):} In ML, datasets often have hundreds of correlated features (standard basis). PCA finds a new orthogonal basis (eigenvectors of the covariance matrix) that points in the directions of maximum variance. Changing coordinates to this new basis decorrelates the data and allows for dimensionality reduction.
    \item \textbf{Quantitative Finance (Risk Factor Modeling):} A portfolio might consist of 500 different stocks (standard basis). Quants map these stocks into a smaller coordinate system of macroeconomic "risk factors" (e.g., interest rates, inflation, market momentum). The change-of-coordinates matrix maps asset returns into factor exposures.
\end{enumerate}

\subsection{B) Model Problem Setup: PCA and Coordinate Transformation}
Let $X \in \mathbb{R}^{n \times m}$ be a centered data matrix representing $m$ samples of $n$-dimensional data (e.g., $n$ financial indicators). The data is currently represented in the standard basis $E$.

We define the sample covariance matrix as:
\[ \Sigma = \frac{1}{m-1} X X^T \]
Since $\Sigma$ is symmetric and positive semi-definite, the Spectral Theorem guarantees we can find an orthonormal basis of eigenvectors $B = \{v_1, v_2, \dots, v_n\}$ such that $\Sigma v_i = \lambda_i v_i$. 

We construct the change-of-coordinates matrix $P_B$:
\[ P_B = \begin{bmatrix} v_1 & v_2 & \dots & v_n \end{bmatrix} \]
Because the basis is orthonormal, $P_B$ is an orthogonal matrix, meaning its inverse is simply its transpose: $P_B^{-1} = P_B^T$.

To transform our entire dataset into the new "Principal Component" coordinate system, we apply the coordinate extraction mapping to every column of $X$:
\[ [X]_B = P_B^{-1} X = P_B^T X \]
In this new coordinate system, the covariance matrix of $[X]_B$ is strictly diagonal. The complex, correlated data in the standard basis has been transformed into completely uncorrelated features in the $B$-basis, drastically simplifying downstream machine learning tasks like regression or clustering.

\newpage
\section{Part 7: Common Variations and Untested Concepts}

\subsection{Change of Basis Between Two Non-Standard Bases}
The homework focuses on translating between a basis $B$ and the standard basis $E$. However, a common exam variation asks you to translate directly between two non-standard bases, $B$ and $C$.

\textbf{The Concept:}
To find the change-of-coordinates matrix from $B$ to $C$, denoted $P_{C \leftarrow B}$, we evaluate the $C$-coordinates of each basis vector in $B$:
\[ P_{C \leftarrow B} = \begin{bmatrix} [b_1]_C & [b_2]_C & \dots & [b_n]_C \end{bmatrix} \]
Then, for any vector $x$, $[x]_C = P_{C \leftarrow B} [x]_B$.

\textbf{Worked Example:}
Let $b_1 = \begin{bmatrix} 1 \\ 2 \end{bmatrix}, b_2 = \begin{bmatrix} 3 \\ 4 \end{bmatrix}$ and $c_1 = \begin{bmatrix} 1 \\ 0 \end{bmatrix}, c_2 = \begin{bmatrix} 1 \\ 1 \end{bmatrix}$. Find $P_{C \leftarrow B}$.

\textit{Step 1:} Find $[b_1]_C$. We need $c_1, c_2$ such that $x_1 c_1 + x_2 c_2 = b_1$.
\[ \begin{bmatrix} 1 & 1 \\ 0 & 1 \end{bmatrix} \begin{bmatrix} x_1 \\ x_2 \end{bmatrix} = \begin{bmatrix} 1 \\ 2 \end{bmatrix} \implies x_2 = 2, x_1 + 2 = 1 \implies x_1 = -1 \]
So, $[b_1]_C = \begin{bmatrix} -1 \\ 2 \end{bmatrix}$.

\textit{Step 2:} Find $[b_2]_C$.
\[ \begin{bmatrix} 1 & 1 \\ 0 & 1 \end{bmatrix} \begin{bmatrix} y_1 \\ y_2 \end{bmatrix} = \begin{bmatrix} 3 \\ 4 \end{bmatrix} \implies y_2 = 4, y_1 + 4 = 3 \implies y_1 = -1 \]
So, $[b_2]_C = \begin{bmatrix} -1 \\ 4 \end{bmatrix}$.

\textit{Step 3:} Construct the matrix.
\[ P_{C \leftarrow B} = \begin{bmatrix} -1 & -1 \\ 2 & 4 \end{bmatrix} \]

\newpage
\section{Part 8: Advanced Diagnostic Testing: "Find the Flaw"}

Below are five problems with subtle, common errors. Your task as the grader is to identify the flaw, explain it, and provide the correct solution.

\subsection{Problem 1}
\textbf{Question:} Given $B = \left\{ \begin{bmatrix} 1 \\ 2 \end{bmatrix}, \begin{bmatrix} 3 \\ 4 \end{bmatrix} \right\}$ and $x = \begin{bmatrix} 5 \\ 6 \end{bmatrix}$, find $[x]_B$.
\textbf{Flawed Solution:} 
To find the coordinates, we multiply the change-of-coordinates matrix by the vector $x$:
\[ [x]_B = P_B x = \begin{bmatrix} 1 & 3 \\ 2 & 4 \end{bmatrix} \begin{bmatrix} 5 \\ 6 \end{bmatrix} = \begin{bmatrix} 5 + 18 \\ 10 + 24 \end{bmatrix} = \begin{bmatrix} 23 \\ 34 \end{bmatrix} \]

\textbf{Correction:}
\textbf{Flaw:} The student multiplied $x$ by $P_B$ instead of $P_B^{-1}$. $P_B$ maps coordinates to vectors, not vectors to coordinates.
\textbf{Correct Step:} $[x]_B = P_B^{-1} x$. 
$P_B^{-1} = \frac{1}{4-6} \begin{bmatrix} 4 & -3 \\ -2 & 1 \end{bmatrix} = \begin{bmatrix} -2 & 1.5 \\ 1 & -0.5 \end{bmatrix}$.
$[x]_B = \begin{bmatrix} -2 & 1.5 \\ 1 & -0.5 \end{bmatrix} \begin{bmatrix} 5 \\ 6 \end{bmatrix} = \begin{bmatrix} -10 + 9 \\ 5 - 3 \end{bmatrix} = \begin{bmatrix} -1 \\ 2 \end{bmatrix}$.

\subsection{Problem 2}
\textbf{Question:} Given $B = \left\{ \begin{bmatrix} 2 \\ 1 \end{bmatrix}, \begin{bmatrix} 5 \\ 3 \end{bmatrix} \right\}$ and $[x]_B = \begin{bmatrix} -1 \\ 2 \end{bmatrix}$, find $x$.
\textbf{Flawed Solution:}
We use the inverse matrix to extract the vector from its coordinates:
\[ P_B = \begin{bmatrix} 2 & 5 \\ 1 & 3 \end{bmatrix} \implies P_B^{-1} = \frac{1}{6-5} \begin{bmatrix} 3 & -5 \\ -1 & 2 \end{bmatrix} = \begin{bmatrix} 3 & -5 \\ -1 & 2 \end{bmatrix} \]
\[ x = P_B^{-1} [x]_B = \begin{bmatrix} 3 & -5 \\ -1 & 2 \end{bmatrix} \begin{bmatrix} -1 \\ 2 \end{bmatrix} = \begin{bmatrix} -3 - 10 \\ 1 + 4 \end{bmatrix} = \begin{bmatrix} -13 \\ 5 \end{bmatrix} \]

\textbf{Correction:}
\textbf{Flaw:} The student used $P_B^{-1}$ to find $x$ from $[x]_B$, but the correct formula is $x = P_B [x]_B$.
\textbf{Correct Step:} $x = \begin{bmatrix} 2 & 5 \\ 1 & 3 \end{bmatrix} \begin{bmatrix} -1 \\ 2 \end{bmatrix} = \begin{bmatrix} -2 + 10 \\ -1 + 6 \end{bmatrix} = \begin{bmatrix} 8 \\ 5 \end{bmatrix}$.

\subsection{Problem 3}
\textbf{Question:} Find the coordinate vector of $p(t) = 4 - 2t + 5t^2$ relative to the standard basis $B = \{1, t, t^2\}$.
\textbf{Flawed Solution:}
We extract the coefficients of the polynomial in decreasing order of degree, as is standard in algebra:
\[ [p]_B = \begin{bmatrix} 5 \\ -2 \\ 4 \end{bmatrix} \]

\textbf{Correction:}
\textbf{Flaw:} The student ordered the coefficients by decreasing degree ($t^2, t, 1$), but the basis $B$ is ordered by increasing degree ($1, t, t^2$). The coordinate vector must strictly match the order of the basis vectors.
\textbf{Correct Step:} $[p]_B = \begin{bmatrix} 4 \\ -2 \\ 5 \end{bmatrix}$.

\subsection{Problem 4}
\textbf{Question:} Find the inverse of the change-of-coordinates matrix $P_B = \begin{bmatrix} 4 & 3 \\ 2 & 1 \end{bmatrix}$.
\textbf{Flawed Solution:}
Using the $2 \times 2$ shortcut, we swap the main diagonal and negate the anti-diagonal:
\[ P_B^{-1} = \begin{bmatrix} 1 & -3 \\ -2 & 4 \end{bmatrix} \]

\textbf{Correction:}
\textbf{Flaw:} The student forgot to multiply the matrix by the scalar $\frac{1}{\det(P_B)}$.
\textbf{Correct Step:} $\det(P_B) = (4)(1) - (3)(2) = 4 - 6 = -2$. 
$P_B^{-1} = -\frac{1}{2} \begin{bmatrix} 1 & -3 \\ -2 & 4 \end{bmatrix} = \begin{bmatrix} -0.5 & 1.5 \\ 1 & -2 \end{bmatrix}$.

\subsection{Problem 5}
\textbf{Question:} Are the polynomials $p_1 = 1+t$, $p_2 = t+t^2$, $p_3 = 1+t^2$ linearly independent?
\textbf{Flawed Solution:}
We map them to $\mathbb{R}^3$ and place them as rows in a matrix to row reduce:
\[ A = \begin{bmatrix} 1 & 1 & 0 \\ 0 & 1 & 1 \\ 1 & 0 & 1 \end{bmatrix} \sim \begin{bmatrix} 1 & 1 & 0 \\ 0 & 1 & 1 \\ 0 & -1 & 1 \end{bmatrix} \sim \begin{bmatrix} 1 & 1 & 0 \\ 0 & 1 & 1 \\ 0 & 0 & 2 \end{bmatrix} \]
Since there are 3 pivots, the rows are linearly independent, so the polynomials are independent.

\textbf{Correction:}
\textbf{Flaw:} While the conclusion is accidentally correct because row rank equals column rank, placing vectors as \textit{rows} does not allow you to find the linear dependence relation $c_1 p_1 + c_2 p_2 + c_3 p_3 = 0$. Coordinate vectors must be placed as \textit{columns}.
\textbf{Correct Step:} Set up the matrix as $A = \begin{bmatrix} [p_1]_B & [p_2]_B & [p_3]_B \end{bmatrix} = \begin{bmatrix} 1 & 0 & 1 \\ 1 & 1 & 0 \\ 0 & 1 & 1 \end{bmatrix}$. Row reducing this column-oriented matrix correctly tests for linear independence by checking for free variables in the equation $Ax = 0$.

\end{document}